\chapter{Transistors: the switch and the amplifier}
\label{ch:transistors}

Everything digital in this book — logic gates, the FPGA fabric you will
program later — is built from one trick, repeated a billion times: a voltage
that controls another current. This chapter introduces the transistor as a
\emph{switch} first (the model you will actually use), then as an
\emph{amplifier} in one honestly-flagged example. When you finish, you will
be able to read a voltage-transfer curve for noise margins, explain why CMOS
logic draws almost no power at rest, and decide whether a given MOSFET can
really drive a given load from a \qty{3.3}{\volt} pin.

% from-zero:begin
\begin{rulebox}[Before you start]
Chapters~2 and~5. \emph{Recap:} a transistor is a switch that another voltage or current controls. Everything digital in this book is made of them.
\end{rulebox}
\begin{wins}[After this chapter you can]
\begin{itemize}
\item explain how a MOSFET works as a voltage-controlled switch;
\item build an inverter and explain why CMOS draws almost no power at rest;
\item read a logic level as ``1'' or ``0'' with noise margins;
\item drive a real load such as a motor or LED with a transistor.
\end{itemize}
\end{wins}
\begin{principle}[How to think about this]
A transistor is a valve: a small signal at the control terminal opens or closes a large path. Use it as a switch and you get digital logic; use it in between and you get amplification. CMOS pairs two complementary valves so that one is always closed, which is why idle chips use almost no power.
\end{principle}
% from-zero:end

\section{The MOSFET: a voltage-controlled switch}
\label{sec:mosfet}

The \emph{MOSFET}\index{MOSFET} (metal-oxide-semiconductor field-effect
transistor) has three terminals we care about — \emph{gate} ($G$),
\emph{drain} ($D$), \emph{source} ($S$) — drawn in
Figure~\ref{fig:mosfet}. The gate sits on an insulating oxide layer: it
controls the device by \emph{electric field}, not by current, so an ideal
gate draws no DC current at all.

\begin{figure}[htbp]
\centering
\begin{tikzpicture}
  % NMOS symbol with terminals labelled and body diode noted
  \node[nmos, scale=1.4](Q) at (2,1.5) {};
  \node[left] at (Q.G) {gate $G$};
  \node[above right] at (Q.D) {drain $D$};
  \node[below right] at (Q.S) {source $S$};
  \draw[->, accent, thick] ($(Q.G)+(-0.6,0)$) -- ($(Q.G)+(-0.1,0)$);
  \node[anchor=west, align=left, font=\small] at (3.6,1.5)
    {$v_{GS} > V_{th}$: channel conducts $D \leftrightarrow S$\\
     gate draws (almost) no DC current};
\end{tikzpicture}

\caption{An NMOS transistor and its three terminals. Notice what the gate
does \emph{not} do: it draws essentially no DC current — it works by field,
not by flow.}
\label{fig:mosfet}
\end{figure}

\begin{definition}[NMOS threshold]
\index{threshold voltage}
An NMOS\index{NMOS} device conducts between drain and source when the
\emph{gate-source voltage} exceeds a device constant $V_{th}$, the
\emph{threshold voltage} (a volt or two for small-signal parts). Below
$V_{th}$ the channel is essentially open.
\end{definition}

The 6.002 pedagogy layers two models on this~\cite{ocw6002}:

\begin{definition}[S and SR models]
\index{MOSFET!S model}\index{MOSFET!SR model}
The \emph{S (switch) model}: $v_{GS} > V_{th}$ $\Rightarrow$ drain and
source are a wire; otherwise an open circuit. The \emph{SR (switch-resistor)
model}: on, the channel behaves as a small resistance $R_{on}$ instead of an
ideal wire — the model you reach for once a load current flows.
\end{definition}

\begin{figure}[htbp]
\centering
\begin{circuitikz}[american]
  % S model: ideal switch
  \node[above] at (1,3.2) {\small\textbf{S model}};
  \draw (1,3) node[circ]{} node[left]{$D$} to[cspst, l_={\scriptsize$v_{GS}>V_{th}$}] (1,1);
  \draw (1,1) node[left]{$S$} node[circ]{};
  \node[align=left, font=\small] at (1,-0.4) {on: wire\\ off: open};
  % SR model: switch + R_on
  \node[above] at (5,3.2) {\small\textbf{SR model}};
  \draw (5,3) node[circ]{} node[left]{$D$} to[cspst, l_={\scriptsize$v_{GS}>V_{th}$}] (5,2)
        to[R, l_=$R_{on}$, -*] (5,0.6) node[left]{$S$} node[circ]{};
  \node[align=left, font=\small] at (5,-0.4) {on: $R_{on}$ (ohms, not zero)\\ off: open};
\end{circuitikz}

\caption{The two switch models. Use S to reason about logic levels; use SR
the moment real current must flow, because $R_{on}$ drops real volts.}
\label{fig:ssr}
\end{figure}

\begin{remark}
Gate current is not literally zero — charging the gate capacitance takes
charge, which is why fast switching burns power (Section~\ref{sec:cmos})
and why FPGAs get warm. "No DC gate current" is a DC statement.
\end{remark}

\section{The NMOS inverter, built by hand}
\label{sec:nmosinv}

Pull the drain up to $V_{DD}$ through a resistor and drive the gate with the
input: when $v_{in} > V_{th}$ the channel conducts and drags the output
down; below threshold the resistor pulls the output up. That is an
\emph{inverter}\index{inverter (logic)} — input 1 gives output 0 and vice
versa. Figure~\ref{fig:vtc} plots its \emph{voltage-transfer
characteristic}\index{voltage-transfer characteristic} (VTC).

\begin{example}[Reading noise margins off a real VTC]\label{ex:vtc}
\textbf{Strategy.} Simulate the NMOS + \qty{10}{\kilo\ohm} pull-up inverter,
then read the four landmarks: the flat output levels $V_{OH}$/$V_{OL}$ and
the input voltages $V_{IL}$/$V_{IH}$ where the curve leaves those flats
(the unit-slope points, read approximately).

\textbf{Solution.} The netlist and the transfer data:
\end{example}

\begin{lstlisting}[caption={\ttfamily code/ch06\_nmos\_inv.cir --- NMOS inverter with a resistive pull-up.},label=lst:nmosinv]
* ch06_nmos_inv.cir -- NMOS + pull-up resistor inverter, .dc transfer
VDD 3 0 dc 3.3
Vin 1 0 dc 0
M1 2 1 0 0 NMOS1 W=10u L=1u
RD 3 2 10k
.model NMOS1 NMOS(vto=1.0 kp=200u lambda=0.02)
.control
dc Vin 0 3.3 0.025
set filetype=ascii
wrdata ch06_nmos_inv.dat v(2)
* static supply current at vin = 3.3 V (inverter holding "low")
dc Vin 3.3 3.3 0.01
print i(VDD)
.endc
.end
\end{lstlisting}

\begin{lstlisting}[language={},caption={Rows of \texttt{ch06\_nmos\_inv.dat}, verbatim --- read $V_{OH}$, $V_{OL}$, and where the knee sits.},label=lst:nmosdata]
vin        vout
0.500      3.3000
1.000      3.3000
1.250      2.6420
1.500      0.7619
1.750      0.2419
2.000      0.1705
3.000      0.0820
3.300      0.0712
\end{lstlisting}

Reading Figure~\ref{fig:vtc}: $V_{OH}=\qty{3.30}{\volt}$,
$V_{OL}\approx\qty{0.07}{\volt}$, $V_{IL}\approx\qty{1.2}{\volt}$ (output
just starting to fall), $V_{IH}\approx\qty{1.7}{\volt}$ (output nearly
bottomed). The margins of Definition~\ref{def:noisemargin} are then
$N M_L \approx 1.2-0.07 \approx \qty{1.1}{\volt}$ and
$N M_H \approx 3.3-1.7 \approx \qty{1.6}{\volt}$.

\begin{figure}[htbp]
\centering
\begin{tikzpicture}
\begin{axis}[width=11cm,height=6.5cm, xlabel={$v_{in}$ (V)}, ylabel={$v_{out}$ (V)},
  xmin=0,xmax=3.3, ymin=0,ymax=3.4, grid=both, legend pos=north east,
  every axis plot/.append style={thick}]
  \addplot[blue] table {figures/ch06/ch06_nmos_inv.dat};
  \addlegendentry{NMOS + $10\,\mathrm{k}\Omega$ pull-up}
  \addplot[red] table {figures/ch06/ch06_cmos_inv.dat};
  \addlegendentry{CMOS}
  % noise-margin landmarks on the NMOS curve
  \draw[dashed, slate] (axis cs:1.2,0) -- (axis cs:1.2,3.3);
  \draw[dashed, slate] (axis cs:1.7,0) -- (axis cs:1.7,3.3);
  \draw[dashed, slate] (axis cs:0,0.071) -- (axis cs:3.3,0.071);
  \node[font=\scriptsize, anchor=west] at (axis cs:0.05,0.45) {$V_{OL}\approx0.07$};
  \node[font=\scriptsize, rotate=90, anchor=south] at (axis cs:1.2,1.0) {$V_{IL}\approx1.2$};
  \node[font=\scriptsize, rotate=90, anchor=south] at (axis cs:1.7,1.0) {$V_{IH}\approx1.7$};
  \node[font=\scriptsize, anchor=west] at (axis cs:0.05,3.15) {$V_{OH}=3.30$};
\end{axis}
\end{tikzpicture}

\caption{Measured inverter transfer curves (both real \texttt{.dc} sweeps).
Notice how much cleaner the CMOS knee is: it snaps over near $V_{DD}/2$ and
actually reaches the rails.}
\label{fig:vtc}
\end{figure}

\section{The static discipline: what makes a level a ``1''}
\label{sec:staticdisc}

Digital only works if stages can be chained: the output of one inverter must
be read correctly by the input of the next, with slack for real-world noise.

\begin{definition}[Static discipline and noise margins]
\label{def:noisemargin}
\index{static discipline}\index{noise margin}
A logic family fixes four voltages: a sender guarantees $v \le V_{OL}$
(``0'') or $v \ge V_{OH}$ (``1''); a receiver accepts $v \le V_{IL}$ as 0
and $v \ge V_{IH}$ as 1, with $V_{OL} < V_{IL} < V_{IH} < V_{OH}$. The gaps
are the \emph{noise margins} $NM_L = V_{IL}-V_{OL}$ and
$NM_H = V_{OH}-V_{IH}$ — how much corruption a wire can absorb before a bit
flips.
\end{definition}

\begin{exercise}[easy]
On Figure~\ref{fig:vtc}'s NMOS curve, is \qty{0.8}{\volt} a safe ``0''
input? Is \qty{1.4}{\volt} a safe ``1'' input? (Use the landmarks just
computed.)
\end{exercise}

\section{CMOS: the pair that draws no power at rest}
\label{sec:cmos}

The resistor pull-up has a defect you can see in the netlist's own printout:
holding the output low means burning $V_{DD}/R$ forever.

\begin{lstlisting}[language={},caption={Last lines of \texttt{ngspice -b ch06\_nmos\_inv.cir} --- the static current while holding output low.}]
i(vdd) = -3.22881e-04
\end{lstlisting}

\qty{0.32}{\milli\ampere} continuously, \emph{per inverter}, for doing
nothing. \emph{CMOS}\index{CMOS} (complementary MOS) replaces the resistor
with a \emph{PMOS}\index{PMOS} transistor — same device, but conducting when
$v_{GS} < V_{th} < 0$, i.e. when the gate is \emph{low}. Now the pull-up is
on exactly when the pull-down is off (Figure~\ref{fig:cmos}): one transistor
of the pair always blocks the rail-to-ground path.

\begin{figure}[htbp]
\centering
\begin{circuitikz}[american, scale=0.9, transform shape]
  % CMOS inverter: PMOS pull-up + NMOS pull-down
  \node[above] at (1.2,4.9) {\small\textbf{inverter}};
  \draw (1.2,4.4) node[pmos, anchor=D](P){};
  \draw (1.2,1.2) node[nmos, anchor=D](N){};
  \draw (P.S) -- (N.D);
  \draw ($(P.S)!.5!(N.D)$) node[circ]{} to[short,-o] ++(1.4,0) node[right]{$Y=\overline{A}$};
  \draw (P.G) -- ++(-0.6,0) |- ($(P.G)!.5!(N.G)+(-0.6,0)$) -- (N.G) ++(-0.6,0);
  \draw ($(P.G)!.5!(N.G)+(-0.6,0)$) node[circ]{} to[short,-o] ++(-0.8,0) node[left]{$A$};
  \draw (N.S) node[ground]{};
  \draw (P.D) ++(0,0.4) node[vdd]{$V_{DD}$} -- (P.D);
  % CMOS NAND: PMOS pair in parallel on top, NMOS pair in series below
  \node[above] at (7.6,4.9) {\small\textbf{NAND}};
  \draw (6.4,4.0) node[pmos, anchor=D](P1){};
  \draw (8.6,4.0) node[pmos, anchor=D](P2){};
  \draw (P1.D) -- ++(0,0.5) -| (P2.D);
  \draw (P1.S) -- (P2.S);
  \draw (7.4,3.6) node[nmos, anchor=D](N1){};
  \draw (7.4,1.6) node[nmos, anchor=D](N2){};
  \draw (N1.S) -- (N2.D);
  \draw (P1.S) -- (N1.D);
  \draw ($(P1.S)!.5!(P2.S)$) node[circ]{} to[short,-o] ++(0.8,0) node[right]{$Y=\overline{A\cdot B}$};
  \draw (N2.S) node[ground]{};
  \draw ($(P1.D)+(0,0.5)$) to[short,-*] ++(0,0.2) node[vdd]{$V_{DD}$};
  \draw (P1.G) to[short,-o] ++(-0.9,0) node[left]{$A$};
  \draw (P2.G) to[short,-o] ++(0.9,0) node[right]{$B$};
  \draw (N1.G) to[short,-o] ++(-0.9,0) node[left]{$A$};
  \draw (N2.G) to[short,-o] ++(-0.9,0) node[left]{$B$};
\end{circuitikz}

\caption{The CMOS inverter (left) and a 2-input CMOS NAND (right): NMOS
network pulls down, PMOS network pulls up, and exactly one network conducts
for any input. Notice the NAND structure — series NMOS (both inputs needed
to pull down), parallel PMOS (either input high pulls up).}
\label{fig:cmos}
\end{figure}

\begin{example}[CMOS static power, measured]\label{ex:cmospower}
\textbf{Strategy.} Sweep the CMOS inverter and read the supply current at
the two resting points — the claim is that both are ~zero, with current only
\emph{during} transitions.

\textbf{Solution.} \texttt{code/ch06\_cmos\_inv.cir} prints
\texttt{i(VDD)} at $v_{in}=0$ and $v_{in}=\qty{3.3}{\volt}$:
\end{example}

\begin{lstlisting}[language={},caption={Static points of the CMOS sweep --- compare with the NMOS inverter's \qty{0.32}{\milli\ampere}.}]
i(vdd) = -3.31000e-12      (vin = 0)
i(vdd) = -3.31000e-12      (vin = 3.3)
\end{lstlisting}

Three picoamps — numerical noise, not real current. Meanwhile the mid-sweep
peak is $\qty{0.33}{\milli\ampere}$ at $v_{in}\approx\qty{1.6}{\volt}$: the
brief window where \emph{both} devices are part-on. That is the whole power
story of modern silicon: \textbf{CMOS burns energy only while
switching}\index{CMOS!power}, which is why processor power scales with
clock frequency.

\begin{principle}[Why CMOS won]
The complementary pair gets rail-to-rail outputs (0 and $V_{DD}$ exactly,
see Figure~\ref{fig:vtc}), a steep transition centred near $V_{DD}/2$, and
zero static current. Resistive pull-ups lose on all three.
\end{principle}

NOR gates swap the networks (parallel NMOS below, series PMOS above).
Every logic function you will meet is these two dual networks — a fact the
FPGA chapters lean on hard.

\begin{exercise}[medium]
Sketch the transistor-level structure of a 2-input CMOS NOR: which devices
are in series and which in parallel, and why must it be the dual of NAND?
\end{exercise}

\section{The BJT: the current-controlled cousin}
\label{sec:bjt}

The \emph{bipolar junction transistor}\index{BJT} (NPN/PNP) controls
collector current with a \emph{base current}: $i_C = \beta\, i_B$, with
$\beta$ of order 100. The mental model shift matters: a MOSFET answers
voltage, a BJT answers \emph{current} — you must supply base current
continuously, which is why discrete switch circuits prefer MOSFETs. BJTs
remain excellent small-signal amplifiers and live inside every op-amp; this
book treats them conceptually only (full biasing design is AoE territory,
flagged honestly).

\section{The transistor as amplifier}
\label{sec:amplifier}

Between "fully on" and "fully off" lies the region the VTC shows as a steep
slope: small changes in $v_{GS}$ produce large changes in $v_{DS}$ — gain.
Bias a single transistor mid-slope, feed it a small signal, and the same
device amplifies. We do exactly one example and then point onward.

\begin{example}[The inverter is also a (bad) amplifier]\label{ex:amp}
\textbf{Strategy.} The VTC slope $-dv_{out}/dv_{in}$ \emph{is} the small-signal
voltage gain. Estimate it from Figure~\ref{fig:vtc}'s steepest point.

\textbf{Solution.} The NMOS curve falls from $\approx\qty{2.6}{\volt}$ to
$\approx\qty{0.76}{\volt}$ between $v_{in}=\qty{1.25}{\volt}$ and
\qty{1.5}{\volt}: a local gain of about $-7$. Biased at
$v_{in}=\qty{1.4}{\volt}$, this inverter is a genuine inverting amplifier —
a terrible one (the bias point wanders with temperature and the pull-up
wastes power), but the same small-signal idea underlies every real
amplifier stage. The art of making this robust — feedback, biasing
networks, differential pairs — is where this book hands you to
AoE~\cite{aacvol3}.
\end{example}

\begin{exercise}[easy]
From the CMOS curve in Figure~\ref{fig:vtc}, estimate its maximum slope in
V/V. Why does a steeper transition make a better gate \emph{and} a better
(though twitchier) amplifier?
\end{exercise}

\section{Driving a load: the switch, for real}
\label{sec:switchload}

\begin{example}[A 2N7000-class MOSFET switching a \qty{12}{\volt} LED strip from \qty{3.3}{\volt} logic]\label{ex:gate}
\textbf{Strategy.} Low-side switch (Figure~\ref{fig:lowhigh}): load between
+\qty{12}{\volt} and drain, source grounded, gate to the logic pin. The
question is never "does it conduct" — it is "does it conduct
\emph{enough}" at the $v_{GS}$ you can actually supply. Check $V_{th}$
against your rail, then simulate the gate-drive sweep.

\textbf{Solution.} \texttt{code/ch06\_switch\_load.cir} sweeps the gate of a
$V_{th}=\qty{2.1}{\volt}$ device driving a \qty{600}{\ohm} strip
($\approx\qty{20}{\milli\ampere}$):
\end{example}

\begin{lstlisting}[language={},caption={Gate-drive sweep, verbatim rows --- the transistor turns on only gradually, and \qty{3.3}{\volt} drive leaves it half-on.},label=lst:switchload]
gate        drain
0.500       12.000
2.000       12.000
3.000        8.061
3.300        5.184
5.000        0.479
\end{lstlisting}

At $v_{GS}=\qty{3.3}{\volt}$ the overdrive is only
$3.3-2.1=\qty{1.2}{\volt}$: the channel passes \qty{11.4}{\milli\ampere}
while the load wants 20 — the strip glows at half brightness and the
transistor eats \qty{5.2}{\volt} of drop. At \qty{5}{\volt} drive the drain
sits at \qty{0.48}{\volt}: solidly on. \textbf{The lesson: match the FET's
specified $R_{on}$ to the gate voltage you actually have} — logic-level
parts specify $R_{on}$ at $v_{GS}=\qty{2.5}{\volt}$ or
\qty{1.8}{\volt} for exactly this reason.

\begin{figure}[htbp]
\centering
\begin{circuitikz}[american, scale=0.9, transform shape]
  % low-side: NMOS between load and ground
  \node[above, align=center] at (1.5,3.6) {\small\textbf{low-side} (ours)};
  \draw (0,3.2) node[vcc]{12 V} to[R, l_={load}] (0,1.8) -- (1.5,1.8);
  \draw (1.5,0.6) node[nmos, anchor=S](M){};
  \draw (M.D) -- (1.5,1.8);
  \draw (M.S) node[ground]{};
  \draw (M.G) to[short,-o] ++(-1.2,0) node[left]{3.3 V logic};
  \node[align=left, font=\small, anchor=west] at (3.0,1.2) {load keeps 12 V on top;\\source on ground --- easy gate drive};
  % high-side: switch above the load
  \node[above, align=center] at (8.5,3.6) {\small\textbf{high-side}};
  \draw (7,3.2) node[vcc]{12 V} -- (8.5,3.2);
  \draw (8.5,2.4) node[nmos, anchor=D](H){};
  \draw (8.5,3.2) -- (H.D);
  \draw (H.S) to[R, l_={load}] (8.5,0) node[ground]{};
  \draw (H.G) to[short,-o] ++(-1.2,0) node[left]{drive};
  \node[align=left, font=\small, anchor=west] at (10.0,1.2) {source rides the load ---\\gate must go \emph{above} the rail};
\end{circuitikz}

\caption{Low-side versus high-side switching. Notice why we always reach
for low-side first: with the source on ground, a \qty{3.3}{\volt} pin is a
known $v_{GS}$; high-side, the source rides the load and the gate needs a
voltage \emph{above the rail} (a bootstrap or charge pump — later study).}
\label{fig:lowhigh}
\end{figure}

\begin{checkpoint}
\textbf{Gate ESD.}\index{ESD!gate} The gate is a few nanometres of oxide; a
charged finger punctures it. Handle MOSFETs and CMOS ICs with static care,
and never leave a gate floating — tie it with a pull-down resistor so the
switch cannot turn on from noise. And this is why Chapter~\ref{ch:transients}'s
flyback diode matters here: the transistor switching a relay coil is
exactly the device the spike kills.
\end{checkpoint}

\begin{rulebox}[High-current switching]
Breadboards and jumper wires are rated for $\sim$\qty{1}{\ampere}. We
simulate anything beyond that — the physics of a \qty{20}{\ampere} switch is
identical, the burns are not.
\end{rulebox}

\section{Logic families and voltage levels}
\label{sec:levels}

Legacy TTL and \qty{5}{\volt} CMOS still exist, but this book lives at
\qty{3.3}{\volt}: modern FPGAs, microcontrollers and the Tang Nano all run
it, and mixing families costs you level shifters\index{level shifter} and
noise-margin headaches. The VTC numbers of Figure~\ref{fig:vtc} are why a
\qty{3.3}{\volt} input is \emph{not} automatically a valid 1 at a
\qty{5}{\volt} input — read the destination's $V_{IH}$ before connecting.

\begin{figure}[htbp]
\centering
\begin{tikzpicture}[font=\small]
  \draw[rounded corners, fill=slate!10] (0,0) rectangle (3,2.4);
  \node[align=center] at (1.5,1.2) {\textbf{analog world}\\ $v(t)$ continuous\\ noise accumulates};
  \draw[rounded corners, fill=accent!12] (5.5,0) rectangle (8.5,2.4);
  \node[align=center] at (7,1.2) {\textbf{digital world}\\ only two levels\\ \texttt{0} / \texttt{1}};
  \draw[->, very thick] (3.2,1.2) -- (5.3,1.2)
    node[midway, above, align=center, font=\small]{static discipline\\ $V_{IL},V_{IH},V_{OL},V_{OH}$\\ + noise margins};
\end{tikzpicture}

\caption{The bridge this chapter builds: analog voltage in, two clean
levels out — guaranteed by the static discipline, not by hope.}
\label{fig:bridge}
\end{figure}

\section{Common mistakes}
\label{sec:transistor-mistakes}

\begin{tripwire}[{Treating $V_{th}$ as a sharp on/off}]
Conduction is gradual (the switch-load sweep rises through 8 V, 5 V on its
way down). "Above threshold" means \emph{starting} to conduct; enough
current for your load is a separate question.
\end{tripwire}

\begin{tripwire}[{Swapping source and drain}]
A discrete NMOS's body diode\index{body diode} points source$\to$drain:
wired backwards, it conducts through the diode no matter what the gate
says. Power packages put a visible diode in the datasheet drawing — look
for it.
\end{tripwire}

\begin{tripwire}[{Leaving the gate floating}]
An un-driven gate drifts on leakage and switches at random — the ESD/gate
checkpoint's pull-down resistor is not decoration.
\end{tripwire}

\begin{tripwire}[{Assuming zero resistance when "on"}]
$R_{on}$ is real: Example~\ref{ex:gate}'s \qty{0.48}{\volt} drain drop at
full drive is \qty{9.6}{\milli\watt} of heat on a small part — check it
before claiming "switch closed".
\end{tripwire}

\section{Summary}
\label{sec:transistor-summary}

\begin{tabularx}{\linewidth}{l X}
\toprule
Concept & Takeaway \\
\midrule
S / SR models & on = wire, or wire + $R_{on}$; off = open \\
VTC & the curve that decides every digital level \\
Noise margins & $NM_L = V_{IL}-V_{OL}$, $NM_H = V_{OH}-V_{IH}$ \\
CMOS & complementary pair, ~0 static power, rail-to-rail \\
Gate drive & $v_{GS}-V_{th}$ sets how \emph{hard} the switch is on \\
Low-side first & source grounded $\Rightarrow$ easy drive \\
Amplifier & same device, biased mid-slope; details live in AoE \\
\bottomrule
\end{tabularx}

\bigskip
Next: flip the switch the other way — same transistors, wired as logic
gates, and the Boolean algebra that drives them.

\begin{exercise}[easy]
A logic family's sender promises $V_{OL}=\qty{0.4}{\volt}$,
$V_{OH}=\qty{2.9}{\volt}$ and the receiver needs $V_{IL}=\qty{0.9}{\volt}$,
$V_{IH}=\qty{2.1}{\volt}$. Compute both noise margins.
\end{exercise}

\begin{exercise}[easy]
Why does an ideal CMOS inverter draw no supply current while holding its
output? Where does the current go during an input transition?
\end{exercise}

\begin{exercise}[easy]
Order these by static power, worst first: CMOS inverter, NMOS + resistor
inverter holding 0, NMOS + resistor inverter holding 1.
\end{exercise}

\begin{exercise}[easy]
Your GPIO pin supplies \qty{3.3}{\volt} and your FET has
$V_{th}=\qty{2.9}{\volt}$. Without simulating: is this a good switch for a
\qty{20}{\milli\ampere} load? What single datasheet number settles it?
\end{exercise}

\begin{exercise}[medium]
Design a 3-input CMOS NAND: draw the pull-down network (series NMOS) and
the pull-up network (parallel PMOS), then state the output expression.
\end{exercise}

\begin{exercise}[medium]
Using Figure~\ref{fig:vtc}'s NMOS curve, estimate $V_{IL}$ and $V_{IH}$ as
the points where $|dv_{out}/dv_{in}|=1$ (the knee is symmetric enough to
eyeball). Recompute both noise margins and compare with the text's
$\approx$1.1/\qty{1.6}{\volt} estimate.
\end{exercise}

\begin{exercise}[medium]
In Example~\ref{ex:gate}, what is the FET's dissipated power at
$v_{GS}=\qty{3.3}{\volt}$ (drain \qty{5.18}{\volt}, current
\qty{11.4}{\milli\ampere})? Compare with the load's and state why the FET
gets hot while the strip stays dim.
\end{exercise}

\begin{exercise}[medium]
Estimate the NMOS inverter's mid-transition supply current from its
resistor: when $v_{out}\approx V_{DD}/2$, what flows through the
\qty{10}{\kilo\ohm} pull-up? Compare with CMOS's measured
\qty{0.33}{\milli\ampere} crossover peak.
\end{exercise}

\begin{exercise}[hard]
Derive the SR-model output low voltage: NMOS on-resistance $R_{on}$ against
pull-up $R_D$ to $V_{DD}$. With $R_{on}=\qty{220}{\ohm}$ and
$R_D=\qty{10}{\kilo\ohm}$ at $V_{DD}=\qty{3.3}{\volt}$, compute $V_{OL}$ and
check it against the simulated \qty{0.071}{\volt}.
\end{exercise}

\begin{exercise}[hard]
A \qty{5}{\volt}-rail legacy gate must read a \qty{3.3}{\volt} CMOS output.
Its $V_{IH}=\qty{3.5}{\volt}$. Quantify the failure, name two fixes, and
state which this book's ecosystem (Tang Nano at 3.3 V logic) avoids the
problem by construction.
\end{exercise}

\section*{Hints and answers}
\begin{description}
\item[Noise margins (easy)] $NM_L=0.9-0.4=\qty{0.5}{\volt}$;
$NM_H=2.9-2.1=\qty{0.8}{\volt}$.
\item[CMOS static (easy)] One device of each pair is off at every level,
so no rail-to-ground path exists; crossover current flows only while both
are part-on, mid-transition.
\item[Static power order (easy)] NMOS holding 0 (resistor conducts),
then CMOS (leakage only); NMOS holding 1 is the odd one: its transistor is
off so it too draws only leakage — the honest order is \{NMOS@0\}
$\gg$ \{NMOS@1, CMOS\}.
\item[Marginal gate (easy)] $v_{GS}-V_{th}=\qty{0.4}{\volt}$ — nearly
closed. $R_{on}$ \emph{at your $v_{GS}$} settles it; a "logic-level" part
specified at \qty{2.5}{\volt} gate is the fix.
\item[3-input NAND (medium)] Three series NMOS below, three parallel PMOS
above: $Y=\overline{A\cdot B\cdot C}$.
\item[Unity-slope margins (medium)] Knee sits near
$v_{in}\approx\qty{1.2}{\volt}$–$\qty{1.7}{\volt}$; margins come out
$\approx$\qty{1.1}{\volt} and $\approx$\qty{1.6}{\volt}, matching the text.
\item[FET dissipation (medium)]
$5.18\times11.4\approx\qty{59}{\milli\watt}$ on the transistor vs
$6.8\times11.4\approx\qty{78}{\milli\watt}$ on the strip
— the FET burns almost half the delivered power.
\item[Pull-up current (medium)] $1.65/10\approx\qty{0.17}{\milli\ampere}$
(\qty{1.65}{\volt} across \qty{10}{\kilo\ohm}); CMOS peaks higher because
its crossover current is channel-limited, not resistor-limited — but only
for the nanoseconds of the transition.
\item[$V_{OL}$ derivation (hard)]
$V_{OL}=V_{DD}\dfrac{R_{on}}{R_{on}+R_D}
=3.3\cdot\dfrac{220}{10220}\approx\qty{0.071}{\volt}$ — dead on the sim.
\item[5 V reading 3.3 V (hard)] Guaranteed-1 input needs
$v\ge\qty{3.5}{\volt}$ but the sender promises only \qty{3.3}{\volt}: zero
margin, guaranteed failure with tolerance drift. Fixes: level-shifter /
gate with lower $V_{IH}$ / run everything at 3.3 V (our choice).
\end{description}

\section*{Notes}
The S/SR model ordering and the static-discipline treatment follow MIT
6.002's digital lectures~\cite{ocw6002}; device physics and the BJT sketch
draw on All About Circuits vol.~III~\cite{aacvol3}; netlists verified on
ngspice-47~\cite{ngspice}. All simulated numbers in this chapter are
verbatim tool output.
